\section{Уроки 1--2: первое подключение, арминг, взлёт и посадка}

\subsection{Цели уроков}
\begin{itemize}
\item Запускать базовый цикл управления без аварийного завершения.
\item Понимать разницу между \texttt{disarm()} и \texttt{land()}.
\item Формировать безопасный блок \texttt{try/finally}.
\end{itemize}

\subsection{Ключевые команды}
\begin{itemize}
\item \texttt{Pioneer()}
\item \texttt{arm()}, \texttt{takeoff()}, \texttt{land()}, \texttt{disarm()}
\item \texttt{close\_connection()}
\end{itemize}

\subsection{Опорные листинги}
\lstinputlisting[style=GeoskanStyle,language=pythoncustom,caption={Минимальное управление моторами и просмотром видеопотока}]{../pioneer-python-example/test.py}
\lstinputlisting[style=GeoskanStyle,language=pythoncustom,caption={Взлёт, выход на точку и посадка}]{../pioneer-python-example/mission_examples/PiZero_test.py}

\subsection{Критерии успешности}
\begin{enumerate}
\item Дрон выполняет арминг и взлёт без ручного вмешательства.
\item После команды посадки двигатели корректно останавливаются.
\item При \texttt{KeyboardInterrupt} аварийный сценарий не оставляет дрон в воздухе.
\end{enumerate}

\begin{taskbox}[Минимальный набор упражнений]
\begin{enumerate}
\item Запустите \texttt{test.py}, проверьте команды \texttt{1}, \texttt{2}, \texttt{ESC}.
\item Запустите \texttt{PiZero\_test.py} и зафиксируйте в журнале время от \texttt{takeoff()} до \texttt{land()}.
\item Добавьте паузу перед \texttt{disarm()} и объясните, зачем она нужна.
\end{enumerate}
\end{taskbox}

\subsection{Контрольные вопросы}
\begin{enumerate}
\item Почему \texttt{land()} предпочтительнее \texttt{disarm()} при штатном завершении?
\item Что должно находиться в блоке \texttt{finally}?
\item Какие признаки устойчивого соединения вы видите в ходе запуска?
\end{enumerate}
